\itemheading{5}{โมเดลการจัดการความรู้ที่ใช้ดำเนินการ}
เสนอใช้ \textbf{KU-KM Canvas Model ร่วมกับวงจร PDCA} เป็นกรอบอธิบายและวางแผนการจัดการความรู้ โดย Canvas ระบุคุณค่า กลุ่มเป้าหมาย ความรู้ ทรัพยากร และการเผยแพร่ ส่วน PDCA กำหนดการทดลอง ทวนสอบ และปรับปรุง กรอบนี้เป็นการสังเคราะห์สำหรับเอกสารเสนอผลงาน ยังไม่ใช่หลักฐานว่าทีมได้ดำเนินกิจกรรม KM ครบทุกขั้นแล้ว

\noindent
\begin{tabularx}{\linewidth}{@{}p{43mm}X@{}}
  \toprule\rowcolor{tableheader}องค์ประกอบ Canvas & การประยุกต์กับระบบลานิสิต\\\midrule
  คุณค่าที่เสนอ (VP) & เชื่อมงานใบลา ข้อมูลสนับสนุน และรายงานให้ตรวจสอบและถ่ายทอดได้\\
  กลุ่มเป้าหมาย (C/S) & นิสิต ที่ปรึกษา ผู้สอน หัวหน้าภาค ผู้บริหาร และเจ้าหน้าที่\\
  ความรู้หลัก (KK) & กฎการลา การนับวัน ลำดับพิจารณา สิทธิ์ ความสัมพันธ์รายวิชา และนิยามตัวชี้วัด\\
  แหล่งความรู้ (SoK) & ซอร์สโค้ด เอกสารโครงการ และระเบียบหรือผู้รับผิดชอบที่ต้องตรวจรับรองเพิ่มเติม\\
  ทรัพยากรหลัก (KR) & ทีมพัฒนา ผู้ดูแลข้อมูล ระบบยืนยันตัวตน ฐานข้อมูล และสภาพแวดล้อมทดสอบ\\
  กิจกรรมหลัก (KA) & ตรวจความต้องการ ออกแบบ พัฒนา ทวนสอบ ถอดบทเรียน และปรับคู่มือ\\
  ตัวชี้วัดสำคัญ (KMe) & ความถูกต้องของรายการและสิทธิ์ ความครบถ้วนข้อมูล เวลาใช้ทำงาน และความเข้าใจของผู้ใช้\\
  ประโยชน์ที่ได้รับ (U) & มีแนวทางกลางในการปฏิบัติงานและติดตามข้อมูล; ผลลดภาระงานต้องประเมินจริง\\
  สินทรัพย์ความรู้ (As) & นิยามข้อมูล แบบจำลอง ระบบ คู่มือ กรณีทดสอบ และบันทึกบทเรียน\\
  คู่ความร่วมมือหลัก (KP) & ภาควิชา ผู้ดูแลข้อมูล ผู้รับผิดชอบระบบตัวตน และผู้แทนแต่ละบทบาทที่จะเชิญทวนสอบ\\
  การบำรุงรักษา (M) & ทบทวนระเบียบ รายชื่อผู้รับผิดชอบ ข้อมูลรายเทอม การสำรอง และสิทธิ์เมื่อมีการเปลี่ยนแปลง\\
  วงจรการเรียนรู้ (LC) & รวบรวมปัญหา ทดลอง ตรวจผล ถอดบทเรียน ปรับปรุง และประเมินซ้ำ\\
  แหล่งจัดเก็บและเผยแพร่ (Ptt) & repository และคู่มือที่ควบคุมเวอร์ชัน; ช่องทางเผยแพร่ตามสิทธิ์ของหน่วยงาน\\\bottomrule
\end{tabularx}

\noindent\textbf{แนวปฏิบัติที่มุ่งให้เกิด:} ทุกการเปลี่ยนกฎหรือข้อมูลสำคัญต้องมีผู้รับผิดชอบ หลักฐานอ้างอิง กรณีตรวจสอบ และคู่มือที่ปรับตามกัน ก่อนประกาศใช้หรือขยายผล
